\section{Critical information controls the response row}
\label{sec:critical-information}

At the critical parameter, the squared edge factor in
\eqref{eq:covariance} is exactly balanced by the number of vertices
at that distance.  To sum a response row, we must therefore obtain
decay from the factors \(a(h_{v_i\to v_{i-1}})\).  We will show that
the observations in one intact infinite side branch supply a loss of
order \(\sqrt t\) at each path step for small \(t\).  The resulting
expected squared row is of order at most \(t^{-1/2}\), which is
integrable at the initial time.

Until the final row estimate, assume
\(b=d-1\ge2\) and \(\theta=b^{-1/2}\).  Constants in the critical
estimates depend only on the fixed degree \(d\).  We work at a fixed
deterministic time \(t>0\) with the raw posterior versions from
Lemma~\ref{lem:posterior}; their continuity is not needed.  For an
oriented edge \(v\to w\), put
\[
 Y=h_{v\to w}(t),\qquad
 M(t)=\E\tanh^2Y,\qquad
 \omega(t)=\min\{\sqrt t,1\}.
\]
The law of \(Y\) is independent of the oriented edge.  It is symmetric:
reversing all spins and all observations preserves their joint law and
negates the conditional field almost surely at this fixed time.
We first determine the scale
of \(M(t)\).  Its lower bound gives path decay, while the whole-site
upper bound will control the early-time energy in the filtering limit.

\subsection{Combining symmetric binary signals}

Combining two branches loses a definite amount of second-moment
information compared with adding their separate contributions.  This
quadratic loss is related to Ioffe's recursive second-moment estimates
for Ising trees \cite[Proposition~4]{Ioffe1996GeneralTrees}.  The following elementary
formulation will allow us to keep the loss explicit.

\begin{lemma}[Symmetric binary-signal fusion]
\label{lem:binary-fusion}
Let \(\varepsilon\) be a uniform sign and let two random elements of
standard Borel spaces be independent conditional on \(\varepsilon\).
Let \(U,V\) be the posterior means of
\(\varepsilon\) from the respective signals.  Suppose that their
marginal laws are symmetric and \(|U|,|V|<1\) almost surely.  If
\(\E_{\mathrm{prod}}\) denotes expectation under the product of these
two marginal laws, then the posterior mean from both signals is
\((U+V)/(1+UV)\), and its second moment is
\begin{equation}
\label{eq:binary-fusion}
 \E U^2+\E V^2-
 \E_{\mathrm{prod}}
   \frac{U^2V^2(2-U^2-V^2)}{1-U^2V^2}.
\end{equation}
In particular this second moment is at most \(\E U^2+\E V^2\).
If \(|U|,|V|\le\theta\), the subtracted term is at least
\(2(1-\theta^2)\E U^2\,\E V^2\).
\end{lemma}

\begin{proof}
For a signal with marginal law \(\nu\) and posterior mean \(u\), the
conditional signal law given \(\varepsilon=\eta\) has density
\(1+\eta u\) with respect to \(\nu\), for \(\eta\in\{-1,1\}\).
Indeed, condition the indicator of \(\{\varepsilon=\eta\}\) on the
signal and use its posterior probability \((1+\eta u)/2\).
Conditional independence then gives joint marginal density \(1+UV\)
relative to the product of the two signal laws.  The signed difference
of the two conditional densities gives the combined posterior mean
\((U+V)/(1+UV)\).  Pushing the product signal law forward to the
product of the marginal laws of \(U,V\), its second moment is
\[
 \E_{\mathrm{prod}}\frac{(U+V)^2}{1+UV}.
\]
The integrand is at most \(2\), so it is legitimate to average it
under the symmetry \(V\mapsto -V\).  The result is
\[
 \frac{U^2+V^2-2U^2V^2}{1-U^2V^2}
   =U^2+V^2-
       \frac{U^2V^2(2-U^2-V^2)}{1-U^2V^2},
\]
which proves \eqref{eq:binary-fusion}.  The loss is nonnegative.
Under the additional bound by \(\theta\), its numerator factor
\(2-U^2-V^2\) is at least \(2(1-\theta^2)\) and its denominator is
at most one.  The expectation of \(U^2V^2\) under the product law is
\(\E U^2\,\E V^2\), proving the final assertion.
\end{proof}

The subadditivity in this lemma can be iterated over the branches of a
vertex.  A collection of already combined signals remains independent
of every unused signal conditional on the spin.  Its posterior mean
remains symmetric, and remains strictly inside \((-1,1)\), since
\[
 1-\left(\frac{u+v}{1+uv}\right)^2
   =\frac{(1-u^2)(1-v^2)}{(1+uv)^2}>0
 \qquad (|u|,|v|<1).
\]
Only the first two child signals below will require the sharper
loss for means bounded by \(\theta\).

\subsection{The amount of information in one component}

For a lower bound, we use a weighted linear estimator in the tradition
of second-moment reconstruction estimates on trees
\cite{EvansKenyonPeresSchulman2000}; see also
Mossel--Peres~\cite[Section~5, Lemma~5.3]{MosselPeres2003}. Here it uses the
Gaussian observations at all depths up to a cutoff.  The cutoff balances
the independent noise against the correlations among the spins.
For the upper bound, the fusion loss gives the complementary estimate.

\begin{lemma}[Critical posterior information]
\label{lem:critical-information}
There are constants \(c,C>0\), depending only on \(d\), such that for
every \(t>0\) and every vertex \(v\),
\begin{equation}
\label{eq:critical-information}
 c\omega(t)\le M(t)\le C\omega(t),
 \qquad
 \E m_v(t)^2\le C\omega(t).
\end{equation}
One may take \(c=1/5\).  In particular both upper bounds are
\(C\sqrt t\) for \(0<t\le1\).
\end{lemma}

\begin{proof}
Root \(T_{v\to w}\) at \(v\).  It is a rooted \(b\)-ary tree; write
\(|u|\) for the depth of \(u\) in it.  For an integer \(n\ge0\), form
the component-observable statistic
\[
 Z_n=\sum_{|u|\le n}\theta^{|u|}X_u(t).
\]
The broadcast covariance is
\(\E[\sigma_u\sigma_z]=\theta^{\dist(u,z)}\).  Since
\(b\theta^2=1\),
\begin{equation}
\label{eq:linear-signal-correlation}
 \E[\sigma_vZ_n]
   =t\sum_{k=0}^n b^k\theta^{2k}=t(n+1).
\end{equation}
The Brownian part of \(Z_n\) has variance \(t(n+1)\) and is independent
of its spin part.

We count the latter part's second moment.  If an ordered pair
\((u,z)\) has depths \(k,\ell\) and its last common ancestor has
depth \(j\), its weighted spin covariance, before the factor \(t^2\),
is
\[
 \theta^{k+\ell+\dist(u,z)}
   =\theta^{2(k+\ell-j)}
   =b^{-(k+\ell-j)}.
\]
When \(j<\min(k,\ell)\), choose the common ancestor, the two distinct
first children in their ordered roles, and then the remaining
descendants.  The exact count is
\[
 b^j\,b(b-1)\,b^{k-j-1}b^{\ell-j-1}
       =(b-1)b^{k+\ell-j-1}.
\]
When \(j=\min(k,\ell)\), the shallower vertex is an ancestor of the
deeper one.  There are \(b^{\max(k,\ell)}\) such ordered pairs;
this includes the \(b^k\) diagonal pairs when \(k=\ell=j\).
Each count is at most \(b^{k+\ell-j}\).  Thus the weighted
contribution for each triple \((k,\ell,j)\) is at most one, and
\begin{align}
 \E Z_n^2
 &\le t(n+1)+t^2
       \sum_{k,\ell=0}^n\bigl(\min(k,\ell)+1\bigr)\notag\\
 &=t(n+1)+t^2\sum_{j=0}^n(n-j+1)^2
 \le t(n+1)+t^2(n+1)^3.
 \label{eq:linear-signal-second-moment}
\end{align}
Because \(Z_n\) is measurable from the component observations,
\(\E[\sigma_vZ_n]=\E[(\tanh Y)Z_n]\).  Cauchy--Schwarz and
\eqref{eq:linear-signal-correlation}--\eqref{eq:linear-signal-second-moment}
therefore give
\begin{equation}
\label{eq:critical-information-lower}
 M(t)\ge\frac{t(n+1)}{1+t(n+1)^2}.
\end{equation}
For \(0<t\le1\), take \(n=\lfloor t^{-1/2}\rfloor\).  Then
\(t^{-1/2}\le n+1\le2t^{-1/2}\), so the last bound is at least
\(\sqrt t/5\).  For \(t\ge1\), take \(n=0\) to obtain
\(M(t)\ge t/(1+t)\ge1/2\).  This proves the lower bound with
\(c=1/5\).

For the upper bound, condition on the component root spin.
Its own observation history and the \(b\) child-component histories
are conditionally independent.  The own signal has posterior mean
\(\tanh X_v(t)\), by the one-spin version of
\eqref{eq:finite-posterior}, and
\[
 q_0(t):=\E\tanh^2X_v(t)\le\E X_v(t)^2=t+t^2.
\]
One child's entire observation history has posterior mean
\(\theta\tanh Y\) about the root spin: conditional on the child spin,
the root has mean \(\theta\) times that spin.  Its second moment is
\(\theta^2M(t)\), and its magnitude is at most \(\theta\).
These statements concern full infinite component histories; their
conditional independence follows from the broadcast construction and
the independent Brownian noises, first for cylinder events and then
for the sigma-fields they generate.

There are at least two child signals because \(b\ge2\).  Combine two
of them first.  Lemma~\ref{lem:binary-fusion} subtracts at least
\(2(1-\theta^2)\theta^4M(t)^2\).  Add the other \(b-2\) child
signals and the own signal by subadditivity.  The combined signal is
the full component history, and hence
\[
 M(t)\le b\theta^2M(t)
       -2(1-\theta^2)\theta^4M(t)^2+q_0(t).
\]
The linear terms cancel at \(b\theta^2=1\).  Since \(q_0(t)\le2t\)
for \(t\le1\), this yields \(M(t)\le C\sqrt t\).
At a whole vertex there are \(d\) child components, so subadditivity
alone gives
\[
 \E m_v(t)^2\le q_0(t)+d\theta^2M(t)\le C\sqrt t
 \qquad(0<t\le1).
\]
Both second moments are at most one for \(t\ge1\), completing
\eqref{eq:critical-information}.
\end{proof}

We have now identified the critical information scale.  To turn its
lower bound into decay along a path, we need the conditional variation
of a side branch's outgoing message, even when the attaching spin has
been specified.

\begin{lemma}[Conditioning a component on a neighboring spin]
\label{lem:component-spin-tilts}
For every integrable function \(f(Y)\) and
\(\eta\in\{-1,1\}\),
\begin{equation}
\label{eq:component-field-tilts}
 \begin{aligned}
 \E[f(Y)\mid\sigma_v=\eta]
   &=\E[f(Y)(1+\eta\tanh Y)],\\
 \E[f(Y)\mid\sigma_w=\eta]
   &=\E[f(Y)(1+\eta\theta\tanh Y)].
 \end{aligned}
\end{equation}
In particular either conditioning leaves the expectation of every
integrable even function of \(Y\) unchanged.
\end{lemma}

\begin{proof}
Let \(\mathcal G=\mathcal F_t^{T_{v\to w}}\).  Since \(\sigma_v\)
is a uniform sign and \(\E[\sigma_v\mid\mathcal G]=\tanh Y\),
conditioning the indicator of \(\{\sigma_v=\eta\}\) on
\(\mathcal G\) proves the first identity.  The component observations
are independent of \(\sigma_w\) conditional on \(\sigma_v\), and
\(\E[\sigma_w\mid\sigma_v]=\theta\sigma_v\).  It follows that
\(\E[\sigma_w\mid\mathcal G]=\theta\tanh Y\); the same indicator
calculation proves the second identity.  The multiplying densities
are bounded, so the formulas apply to every marginally integrable
\(f(Y)\).  If \(f\) is even, its product with \(\tanh Y\) is odd.
Symmetry of the marginal law of \(Y\) makes that correction vanish.
\end{proof}

\subsection{Attenuation from an intact side branch}

The next argument uses the full infinite component beyond one side
edge.  Its conditioning specifies the spins on a deterministic path,
while leaving the selected side branches unobserved.  This is the
setting in which the information lower bound supplies a loss uniformly
over the contribution from the rest of the tree.

\begin{lemma}[Conditional path attenuation]
\label{lem:path-attenuation}
There is \(c_0>0\), depending only on \(d\), such that for every
\(t>0\), every \(z\in\mathbb R\), and \(\eta\in\{-1,1\}\),
\begin{equation}
\label{eq:shifted-attenuation}
 \E\left[
   a\bigl(z+\psi(h_{v\to w}(t))\bigr)^2
       \,\middle|\,\sigma_w=\eta\right]
 \le1-c_0\omega(t).
\end{equation}
For every deterministic simple path \(v_0,\ldots,v_r\), \(r\ge1\),
\begin{equation}
\label{eq:conditional-path-attenuation}
 \E\left[
   \prod_{i=1}^r a(h_{v_i\to v_{i-1}}(t))^2
       \,\middle|\,\sigma_{v_0},\ldots,\sigma_{v_r}\right]
 \le e^{-c_0\omega(t)r}.
\end{equation}
The same bound holds for the reverse product
\(\prod_{i=1}^r a(h_{v_{i-1}\to v_i}(t))^2\), and both bounds
also hold without conditioning on the path spins.
\end{lemma}

\begin{proof}
Set \(G=\psi(Y)\) and condition on the attaching spin
\(\sigma_w=\eta\).  We have \(|G|\le\beta\).  The second identity
in \eqref{eq:component-field-tilts}, symmetry, and Cauchy--Schwarz give
\[
 \E[G^2\mid\sigma_w=\eta]=\E G^2,\qquad
 |\E[G\mid\sigma_w=\eta]|^2
   =\theta^2|\E[G\tanh Y]|^2
   \le\theta^2\E G^2.
\]
Moreover \(|\psi(y)|\ge\theta|\tanh y|\), since
\(|\operatorname{atanh}u|\ge|u|\).  Consequently
\begin{equation}
\label{eq:side-message-variance}
 \Var(G\mid\sigma_w=\eta)
   \ge(1-\theta^2)\E G^2
   \ge(1-\theta^2)\theta^2M(t).
\end{equation}
The conditional variance is the useful quantity here: unlike an
uncentered second moment, it cannot be eliminated by adding a shift.

For any random variable \(G\) supported on \([-\beta,\beta]\), there
is a constant \(c_\beta>0\) such that
\begin{equation}
\label{eq:shifted-tanh-variance}
 \E\tanh^2(z+G)\ge c_\beta\Var(G)
 \qquad\text{for every }z\in\mathbb R.
\end{equation}
To see this, first suppose \(|z|\le\beta+1\).  Then
\(|z+G|\le2\beta+1\).  The ratio \(\tanh^2 u/u^2\), extended
continuously to be one at \(u=0\), has a positive minimum on that
compact interval.  Hence
\(\E\tanh^2(z+G)\ge c_\beta\E(z+G)^2\ge c_\beta\Var(G)\).
If \(|z|>\beta+1\), then \(|z+G|>1\), so the left side is at least
\(\tanh^2(1)\); use \(\Var(G)\le\beta^2\) and decrease
\(c_\beta\) if needed.  This proves \eqref{eq:shifted-tanh-variance}
for the conditional laws in \eqref{eq:side-message-variance} as well.

For \(u=\tanh^2h\), the elementary bounds on \(a\) give
\[
 a(h)^2\le a(h)
   =1-\frac{(1-\theta^2)u}{1-\theta^2u}
   \le1-(1-\theta^2)u.
\]
Apply this with \(h=z+G\).  Equations
\eqref{eq:shifted-tanh-variance} and \eqref{eq:side-message-variance},
followed by \(M(t)\ge\omega(t)/5\), prove
\eqref{eq:shifted-attenuation}, for example with
\[
 c_0=\min\left\{\frac12,\,
       \frac{c_\beta(1-\theta^2)^2\theta^2}{5}\right\}.
\]
The same estimate may be used with a random shift after conditioning
on data that fix the shift and are independent of the selected branch
given its attaching spin.

For the path estimate, choose at each target \(v_i\), using only the
path geometry, one neighbor \(s_i\) that is not on the path.  There are
\(d-2\ge1\) choices at an internal target and \(d-1\) at the last
target.  The components \(R_i=T_{s_i\to v_i}\) are disjoint and
infinite.  In particular, when \(d=3\) an internal target still has
one full infinite side component.

Let
\[
 \mathcal H=\sigma\bigl(
   \sigma_{v_0},\ldots,\sigma_{v_r},\
   X_u(s):u\notin\textstyle\bigcup_iR_i,\ 0\le s\le t\bigr).
\]
Conditional on \(\mathcal H\), the selected
component observation arrays are independent, and the law in \(R_i\)
depends only on its attaching spin \(\sigma_{v_i}\).  This follows
directly by rooting the broadcast at \(v_0\): the multiplier on each
attaching edge, the multipliers inside its selected component, and its
Brownian motions form one independent group.  The outside observations
use none of these groups once the path spins are specified.
The assertion for the full
observation arrays follows from the corresponding assertion on finite
cylinder events.

At this fixed time, the cavity identity gives, almost surely,
\[
 h_{v_i\to v_{i-1}}
   =\psi(h_{s_i\to v_i})+
      \left(X_{v_i}
        +\sum_{\substack{u\sim v_i\\u\ne v_{i-1},\,s_i}}
                      \psi(h_{u\to v_i})\right).
\]
When all other selected arrays are held fixed, the parenthesized
quantity is a fixed shift independent of the array in \(R_i\)
conditional on \(\sigma_{v_i}\).  Integrate \(R_1\) first.
Every factor at a later target \(v_j\), \(j>1\), uses its component
on the side excluding \(v_{j-1}\), and hence does not use \(R_1\).
Thus \eqref{eq:shifted-attenuation} removes the first squared factor
at a cost at most \(1-c_0\omega(t)\).  Repeating in the order
\(R_2,\ldots,R_r\) gives
\[
 \E\left[\prod_{i=1}^r a(h_{v_i\to v_{i-1}})^2
       \,\middle|\,\mathcal H\right]
 \le (1-c_0\omega(t))^r
 \le e^{-c_0\omega(t)r}.
\]
Taking a conditional expectation given only the path spins proves
\eqref{eq:conditional-path-attenuation}.  Reversing the deterministic
path gives the stated reverse product.  Finally average over its spins.
\end{proof}

The loss above comes from the posterior information in an intact
infinite branch.  It is not a claim about arbitrarily truncated
branches: a finite boundary will require a separate estimate using
the target's own Gaussian observation.  For the present response row,
the infinite-branch estimate is exactly what is needed.

\subsection{Time-integrable response rows}

Return to the full range \(0<\theta\le b^{-1/2}\).  For the
progressive coefficients in \eqref{eq:covariance}, define the extended
nonnegative sum
\[
 A_v(t)=\sum_{u\in\mathcal V}c_{vu}(t)^2.
\]
It is measurable as a countable sum.  The following estimate is the
input that permits division by the posterior variance when the
filtering equation is converted to a field equation.

\begin{proposition}[Time-integrable response rows]
\label{prop:row-energy}
For every vertex \(v\), there is \(C<\infty\), depending only on \(d\)
at the critical parameter, such that
\begin{equation}
\label{eq:row-energy}
 \begin{aligned}
 A_v(t)&\le1+\frac{d\theta^2}{1-b\theta^2}
      &&\text{if }b\theta^2<1,\quad t\ge0,\\
 \E A_v(t)&\le\frac{C}{\omega(t)}
      &&\text{if }b\theta^2=1,\quad t>0,\\
 \E\int_0^T A_v(t)\dd t&<\infty
      &&\text{for every }T<\infty\text{ in either case}.
 \end{aligned}
\end{equation}
The same fixed-time and integrated upper bounds hold with
\(\sum_u C_{vu}(t)^2\) in place of \(A_v(t)\).
\end{proposition}

\begin{proof}
There are \(db^{r-1}\) vertices at distance \(r\ge1\) from \(v\).
If \(b\theta^2<1\), the pointwise bound \(a\le1\) gives
\[
 A_v(t)\le1+\sum_{r\ge1}db^{r-1}\theta^{2r}
       =1+\frac{d\theta^2}{1-b\theta^2}.
\]
At equality, \(db^{r-1}\theta^{2r}=d/b\).  Tonelli's theorem and
Lemma~\ref{lem:path-attenuation} give, for \(t>0\),
\[
 \E A_v(t)
  \le1+\frac db\sum_{r\ge1}e^{-c_0\omega(t)r}
  =1+\frac{d/b}{e^{c_0\omega(t)}-1}
  \le1+\frac{d}{bc_0\omega(t)}
  \le\frac{C}{\omega(t)}.
\]
Here \(e^x-1\ge x\) and \(\omega(t)\le1\).
Since
\[
 \int_0^T\frac{\dd t}{\omega(t)}
   =
   \begin{cases}
     2\sqrt T,&0\le T\le1,\\
     T+1,&T\ge1,
   \end{cases}
\]
another application of Tonelli proves the integrated bound.
Finally \(\sum_u C_{vu}^2=(1-m_v^2)^2A_v\le A_v\) by
\eqref{eq:covariance}.
\end{proof}

In particular, \(A_v(t)\) is finite almost surely at each fixed
positive time and is finite for almost every time almost surely on
each compact horizon.  At criticality the raw sum at \(t=0\) is
infinite: all fields vanish, \(a(0)=1\), and every positive-distance
shell contributes \(d/b\).  That single time has zero Lebesgue measure.
The integrated estimate, rather than a bound at time zero, is the
conclusion used in the construction of the infinite innovation equation.
